\subsection{通常のPrym多様体との関係}
\label{subsec:classical-prym-context-v27}

同じnetから作る二つの被覆を整理しておく。
対称表示が定義される体上で、$C:ac-b^2=0$ に対応する非分岐二重被覆は
\[
 \widetilde C:\quad a(x)=r^2,\quad b(x)=rs,\quad c(x)=s^2
 \quad\subset\mathbb P^4
\]
という種数五の標準モデルを持つ。
この三つの二次超曲面のnetの判別式は
\[
 \Delta_N\ \cup\ \{4\lambda_0\lambda_2-\lambda_1^2=0\}
\]
という平面五次曲線である。
Bruinの構成\cite[Theorem 5.1, case 4]{BruinPrymGenusThree}では、一般の場合、
二次元Prym多様体は種数二曲線
\[
 F:\quad y^2=-\det(A_0+2tA_1+t^2A_2)
\]
のJacobianとして表される。その六つの分岐点は上の二次曲線と $\Delta_N$ の交点である。
一方、本稿の $E\to\Delta_N$ は、平面内の退化二次曲線の二直線成分を選ぶ
\emph{エタール}被覆である。
種数五の標準モデルの $W^1_4$ を使って判別式上に作る、六点で分岐する被覆とは異なる。
従って $\mathcal R_3\dashrightarrow Y_0(2)$ はnetに付随する補助不変量を記述している。
